\chapter{Introduction to AI-Based Automated Log Analysis}
This chapter introduces the context of emulation-based pre-silicon validation and the role of log analysis in modern System-on-Chip (SoC) development. It highlights the limitations of manual log inspection, surveys recent literature on log parsing and anomaly detection, defines the motivation and the formal problem statement of the project, and lists the objectives. A brief outline of the proposed methodology and the organization of the report is also presented.

\section[Introduction]{\textbf{Introduction}}
In modern semiconductor development, hardware and software are designed together (co-design) to ensure that system functionality and performance targets are met before the silicon is fabricated. Emulation platforms enable testing of system software even before the actual hardware is available, which is critical for catching bugs early and reducing the overall product development cycle. Multiple software test cases are executed on these emulation platforms, and each test generates detailed log files that capture the system behaviour and performance counters of interest.

These logs contain important metrics such as bandwidth, clock cycles, and instructions per second (IPS), along with workload-specific counters for Coproc\_1 (matrix engine), Coproc\_2 (vector engine), and the scalar core. They also contain information about stalls, memory pressure (spill/fill bytes), tiling, scheduling, and build/configuration parameters. These metrics are essential for performance evaluation, regression comparison and root-cause analysis. However, with more than one hundred logs generated per regression run, manual extraction and analysis of metrics becomes impractical.

To overcome this limitation, the proposed system integrates LogPAI (using the Drain algorithm) for structured log parsing and LogAI for machine learning based analysis, along with an LLM-based module that produces high-level summaries and assists in root-cause analysis. The output is delivered as standardized per-regression workbooks, a final multi-option comparison workbook, dashboards and email alerts, so that the same analysis can be reproduced by any engineer in the core team. The high-level data flow of the proposed system starts from raw log generation on the emulation platform, proceeds through structured parsing, metric extraction and analysis, and ends with report generation and notifications. The detailed four-stage methodology is presented in Chapter 3.

\section[Literature Review]{\textbf{Literature Review}}

A literature review is a scholarly survey of the current state of knowledge on a topic, including substantive findings as well as theoretical and methodological contributions. The literature on automated log analysis has expanded rapidly in the recent years, driven by the growing scale of system logs in cloud, telemetry and pre-silicon validation environments. This section surveys five recent papers that directly relate to the techniques used in this project.

\subsection{Sample}
Recent literature on log parsing and anomaly detection focuses on three broad themes: (i) semantic and LLM-driven parsing of unstructured logs, (ii) cross-system generalization of anomaly detection models, and (iii) scalable, low-latency parsing for online use. The summaries below capture the key findings from each paper and indicate how they motivate the design choices in this work.

\subsubsection[Review types]{\textbf{Review types}}

Zhang et al. (2025), in \textit{SemanticLog: Towards Effective and Efficient Large-Scale Semantic Log Parsing}, demonstrate that LLM-based semantic parsing improves both accuracy and efficiency of large-scale log analysis while preserving privacy. The paper motivates the use of semantic templates over purely regex-based parsers when the log corpus contains many heterogeneous formats.

Liu et al. (2025), in \textit{LUK: Empowering Log Understanding With Expert Knowledge From Large Language Models}, show that incorporating LLM-derived expert knowledge significantly enhances downstream log understanding and analysis performance. This justifies the integration of an LLM-based summarization and root-cause analysis stage in the proposed pipeline.

\subsubsection[Process and product]{\textbf{Process and product}}

Wang et al. (2025), in \textit{LogDLR: Unsupervised Cross-System Log Anomaly Detection Through Domain-Invariant Latent Representation}, introduce domain-invariant representations that improve anomaly detection across different systems with varying log formats. The work is relevant because the emulation platform produces logs from several sub-systems with different formats but similar underlying performance semantics.

Chen et al. (2025), in \textit{EPAS: Efficient Online Log Parsing via Asynchronous Scheduling} (IEEE ICDE), demonstrate that asynchronous scheduling of parsing tasks improves the speed and efficiency of log parsing in real-time systems. This is directly applicable to the automated regression workflow, where parsing latency is a key constraint.

\subsubsection{\textbf{Page limitation}}

Li et al. (2025), in \textit{Practitioners' Expectations on Log Anomaly Detection} (IEEE Transactions on Software Engineering), highlight the gap between academic research and real-world industrial requirements in log anomaly detection systems. Their findings reinforce the need for tooling that produces interpretable, actionable summaries, rather than only raw model scores - which is one of the central design goals of this project.

Together, these works provide the technical and practical foundation for the design of the proposed pipeline. They support the choice of the Drain algorithm (LogPAI) for template extraction, the use of LogAI for clustering and anomaly detection on parsed logs, and the integration of an LLM API for natural-language summarization and root-cause direction.

\subsubsection{\textbf{Plagiarism}}

All concepts borrowed from these papers are referenced in the bibliography. Direct quotations from the papers have been avoided; the discussion paraphrases the original ideas while preserving their technical meaning, in line with standard academic practice.

\subsubsection{How to add Reference}
The reference papers are managed using \texttt{Jabref} and cited with the \verb|\cite{}| command. A representative citation, for example to a textbook, would look like this \cite{Razavi2000}. The complete list of references is given at the end of the report.

\section[Motivation]{\textbf{Motivation}}

The motivation for this project comes directly from the daily challenges faced by the pre-silicon performance validation team. The key issues are:

\begin{itemize}
\item Manual analysis of a large number of test logs (typically more than 100 per regression run) is highly repetitive and consumes a significant fraction of an engineer's working day.
\item Engineers spend more effort on extracting data from raw logs than on actual root-cause analysis and problem solving.
\item Without an automated comparison framework, it is difficult to quickly determine whether an observed degradation is caused by hardware (architecture, clock frequency) or by software (compiler version, build flags).
\item Lack of automated summarization delays communication of regression results to managers and to other cross-functional teams.
\item Existing tools individually handle parsing, clustering or visualization, but none of them provide a unified hybrid AI-based pipeline that combines structured parsing, ML-based analysis and LLM-based reporting.
\end{itemize}

There is therefore a strong industrial need for an efficient, faster and reproducible debugging tool that can keep up with aggressive product development timelines.

\section[Problem statement]{\textbf{Problem statement}}

To overcome the slow and error-prone manual analysis of extensive log data generated from large-scale emulation tests, and to effectively identify performance issues and their root cause, an automated system is required for efficient log extraction, analysis and summarization that improves debugging and decision-making across the validation team.

\section[Objectives]{\textbf{Objectives}}
The objectives of the project are:
\begin{enumerate}
\item To replace manual log-grepping with a single automated pipeline that runs per regression and produces repeatable outputs for any engineer in the core team.
\item To compress regression analysis time from several hours to approximately two minutes per regression by automating the extraction and formatting of results.
\item To standardize customer-request reporting through per-regression three-sheet workbooks and a final multi-option comparison workbook (Cascade baseline vs. Nova-* variants), and to enable fast degradation detection and root-cause direction using automatically computed percentage differences and stall/bandwidth indicators such as Coproc\_2 stalls and DMA-bound signals.
\end{enumerate}

\section[Brief Methodology of the project]{\textbf{Brief Methodology of the project}}
The methodology of this project is structured into four stages: data collection, log processing, analysis engine and output \& delivery. In the data collection stage, software test cases are executed on the emulation platform and raw log files are ingested and stored for offline processing. In the log processing stage, raw logs are pre-processed and parsed using LogPAI's Drain algorithm to produce structured log templates and feature vectors. The analysis engine combines LogAI-based machine learning (clustering and anomaly detection), regex-driven metric extraction (cycles, IPS, bandwidth, Coproc\_1/Coproc\_2/scalar counters, stalls) and an LLM-based stage for summary generation and root-cause analysis. The final stage produces CSV / XLSX reports, a dashboard view, and email/status notifications. The detailed methodology is presented in Chapter 3.

\section[Assumptions made / Constraints of the project]{\textbf{Assumptions made / Constraints of the project}}

The system is developed under the following assumptions and constraints:
\begin{itemize}
\item All log files follow the format produced by the in-house emulation platform; cross-vendor log formats are not in scope.
\item The set of performance counters extracted (cycles, IPS, bandwidth, Coproc\_1/Coproc\_2/scalar counters, stalls, spill/fill bytes etc.) is fixed for the current release and matches the counters available on the platform under study.
\item LLM-based summary and root-cause analysis is advisory in nature; final engineering decisions are taken by validation engineers after reviewing the structured workbooks.
\item The pipeline is validated on a fixed test suite of twelve representative graphs and on six silicon configurations (one Cascade baseline plus five Nova variants).
\item Network access to the LLM API is assumed to be available within the development environment.
\end{itemize}

\section[Organization of the report]{\textbf{Organization of the report}}

This report is organized as follows.
\begin{itemize}
\item Chapter 2 discusses the theoretical background, including emulation-based pre-silicon validation, log parsing using the Drain algorithm (LogPAI), LogAI-based ML analysis, and the use of LLMs for log summarization, along with the software tools used in the project.
\item Chapter 3 presents the detailed analysis and design of the proposed system, covering data collection, log processing, the analysis engine, output \& delivery, and the list of performance parameters extracted from each log.
\item Chapter 4 describes the implementation of the pipeline, including the Python-based modules for log ingestion, regex extraction, structured workbook generation, comparison engine and LLM integration.
\item Chapter 5 presents the results obtained, including the degradation summary across all six options, the IPS comparison of Nova-Alpha vs.\ Cascade, and a detailed root-cause analysis of MosaicNet\_W8A8.
\item Chapter 6 concludes the report and outlines possible directions for future work, along with the learning outcomes from the project.
\end{itemize}
